\chapter{\IfLanguageName{dutch}{Stand van zaken}{State of the art}}%
\label{ch:stand-van-zaken}

% Tip: Begin elk hoofdstuk met een paragraaf inleiding die beschrijft hoe
% dit hoofdstuk past binnen het geheel van de bachelorproef. Geef in het
% bijzonder aan wat de link is met het vorige en volgende hoofdstuk.

% Pas na deze inleidende paragraaf komt de eerste sectiehoofding.

%Dit hoofdstuk bevat je literatuurstudie. De inhoud gaat verder op de inleiding, maar zal het onderwerp van de bachelorproef *diepgaand* uitspitten. De bedoeling is dat de lezer na lezing van dit hoofdstuk helemaal op de hoogte is van de huidige stand van zaken (state-of-the-art) in het onderzoeksdomein. Iemand die niet vertrouwd is met het onderwerp, weet nu voldoende om de rest van het verhaal te kunnen volgen, zonder dat die er nog andere informatie moet over opzoeken \autocite{Pollefliet2011}.

%Je verwijst bij elke bewering die je doet, vakterm die je introduceert, enz.\ naar je bronnen. In \LaTeX{} kan dat met het commando \texttt{$\backslash${textcite\{\}}} of \texttt{$\backslash${autocite\{\}}}. Als argument van het commando geef je de ``sleutel'' van een ``record'' in een bibliografische databank in het Bib\LaTeX{}-formaat (een tekstbestand). Als je expliciet naar de auteur verwijst in de zin (narratieve referentie), gebruik je \texttt{$\backslash${}textcite\{\}}. Soms is de auteursnaam niet expliciet een onderdeel van de zin, dan gebruik je \texttt{$\backslash${}autocite\{\}} (referentie tussen haakjes). Dit gebruik je bv.~bij een citaat, of om in het bijschrift van een overgenomen afbeelding, broncode, tabel, enz. te verwijzen naar de bron. In de volgende paragraaf een voorbeeld van elk.

Dit hoofdstuk biedt een analyse van de huidige situatie in de mainframe-industrie,
bestaande competentiemodellen en opleidingsmethoden. Het bouwt voort op de
inleiding door dieper in te gaan op de context en relevantie van mainframes en
z/OS systeembeheer, en legt de basis voor de methodologie en het ontwikkeldeframe
work in de volgende hoofdstukken. Hierbij wordt verwezen naar relevante litera
tuur om de stand van zaken te schetsen.

\section{\IfLanguageName{dutch}{Wat zijn mainframes}{What are mainframes}}
\label{sec:wat-zijn-mainframes}

Vooraleer dieper in te gaan op de kern van deze bachelorproef, is het noodzakelijk om een duidelijk beeld te schetsen van wat mainframes zijn. Mainframes zijn krachtige, hoogperformant computersystemen die ontworpen zijn voor het verwerken van grote hoeveelheden transacties en data met een uitzonderlijk hoge betrouwbaarheid en beschikbaarheid~\autocite{IBMMainframe}.

Mainframes functioneren continu (24/7) en halen beschikbaarheidsniveaus van 99,999\% uptime~\autocite{IBM_zOS_Concepts}, wat neerkomt op slechts enkele minuten ongeplande downtime per jaar. In combinatie met hun schaalbaarheid en ingebouwde fouttolerantie maakt dit hen bijzonder geschikt voor omgevingen waar systeemuitval grote financiële of maatschappelijke gevolgen kan hebben.

Naast mainframes bestaan er andere platformen voor het uitvoeren van toepassingen en het opslaan van gegevens, zoals Linux- en Windows-servers. Hoewel deze systemen flexibel en breed inzetbaar zijn, verschillen zij fundamenteel van mainframes op het vlak van architectuur, geïntegreerde betrouwbaarheid en schaalbaarheid. Waar klassieke servers vaak afhankelijk zijn van clustering en externe virtualisatietechnologieën om vergelijkbare redundantie te realiseren, zijn deze eigenschappen inherent ingebouwd in mainframe-architecturen.
\subsection{\IfLanguageName{dutch}{Waarom IBM Mainframes of z/OS}{Why IBM Mainframes or z/OS}}
\label{sec:waarom-ibm-mainframe-of-zos}
Binnen de huidige mainframemarkt wordt het overgrote deel van de mainframes geleverd door IBM, dat meer dan 90\% van het wereldwijde marktaandeel in handen heeft~\autocite{MarketReportsWorld2026}. Daarnaast vertegenwoordigen IBM Z-systemen ongeveer 63,25\% van de wereldwijde mainframemarkt in 2025~\autocite{MordorIntelligence2026}. Dit onderbouwt de keuze om z/OS als focusplatform te nemen, aangezien IBM Z het dominante commerciële mainframe-ecosysteem vormt.

Alternatieve mainframesystemen, zoals het stack-gebaseerde besturingssysteem MCP van Unisys, hebben een beperkter marktaandeel en een kleinere aanwezigheid in hedendaagse enterprise-omgevingen.

\subsection{\IfLanguageName{dutch}{Hardware-onderscheid}{Hardware difference}}
\label{sec:hardware-onderscheid}

Mainframes onderscheiden zich van andere platformen door hun gespecialiseerde hardware-architectuur, die ontworpen is met het oog op hoge betrouwbaarheid en prestaties. Ze beschikken over redundante componenten, geavanceerde fouttolerantiemechanismen en ondersteunen duizenden gelijktijdige gebruikers en applicaties. Daarnaast zijn mainframes in staat om grote hoeveelheden gegevens te verwerken en complexe transacties efficiënt af te handelen.

Linux- en Windows-servers daarentegen zijn doorgaans opgebouwd uit standaard hardwarecomponenten en worden vaak geconfigureerd voor specifieke toepassingen of workloads. Hoewel dergelijke systemen schaalbaar zijn, bereiken zij doorgaans niet hetzelfde niveau van geïntegreerde betrouwbaarheid en beschikbaarheid als mainframes. Om vergelijkbare schaalbaarheid te realiseren, worden servers vaak gegroepeerd in clusters. Dit kan echter leiden tot bijkomende complexiteit op het vlak van beheer, configuratie en onderhoud~\autocite{EES2025}.

\subsection{\IfLanguageName{dutch}{Virtualisatie, logische partities en resourcebeheer}{Virtualization, logical partitions and resource management}}
\label{sec:virtualisatie-en-resourcebeheer}
Mainframes bieden geavanceerde virtualisatiemogelijkheden, waardoor één fysiek systeem kan worden opgedeeld in meerdere logische partities (LPARs). Deze partities worden beheerd door de ingebouwde hypervisor PR/SM (Processor Resource/System Manager), die zorgt voor een strikte isolatie tussen omgevingen. Elke LPAR kan een eigen besturingssysteem draaien, zoals z/OS, z/VM of Linux on Z, en beschikt over toegewezen resources zoals CPU, geheugen en I/O-capaciteit.\autocite{IBMIntroLPAR}

Een belangrijk kenmerk van deze architectuur is de mogelijkheid tot dynamische resource-allocatie. Resources kunnen tijdens runtime worden aangepast zonder het systeem te herstarten, wat bijdraagt aan flexibiliteit en hoge beschikbaarheid. Daarnaast maakt Workload Management (WLM) het mogelijk om prioriteiten toe te kennen aan verschillende workloads, zodat kritieke transacties voorrang krijgen bij piekbelasting.\autocite{IBM_zOS_Concepts}

Deze combinatie van virtualisatie, isolatie en intelligent resourcebeheer stelt organisaties in staat om meerdere bedrijfskritische toepassingen efficiënt en veilig op één fysiek platform te consolideren, zonder in te boeten aan stabiliteit of prestaties.

\subsection{\IfLanguageName{dutch}{Soorten besturingsystemen in IBM mainframes}{Types of operating systems in IBM mainframes}}
\label{sec:soorten-besturingsystemen-in-ibm-mainframe}
Hieronder vindt u een overzicht van de verschillende besturingssystemen die op IBM-mainframes kunnen draaien en eventueel in een logische partitie (LPAR) worden geïnstalleerd/gedraaid. Elk van deze besturingssystemen heeft zijn eigen kenmerken, toepassingen.
\begin{itemize}
	\item \textbf{z/OS:} Het belangrijkste enterprise-besturingssysteem van IBM Z voor bedrijfskritische workloads, met sterke focus op beschikbaarheid, beveiliging en grootschalige batch- en transactieverwerking \autocite{IBMRedbook_ABCzOS,IBMRedbook_zOSBasics,IBMzOSOverview}.
	
	\item \textbf{z/VM:} IBM’s virtualisatieplatform (hypervisor) dat meerdere virtuele machines op één mainframe mogelijk maakt en vaak wordt gebruikt om z/OS en Linux omgevingen te consolideren \autocite{IBMRedbook_zOSBasics}.
	
	\item \textbf{z/TPF:} Een gespecialiseerd besturingssysteem voor extreem hoge transactiesnelheden en lage latency, typisch in sectoren zoals luchtvaart en betalingsverkeer \autocite{IBMRedbook_zOSBasics,ibm_ztpf_2026}.
	
	\item \textbf{Linux on Z:} Linuxdistributies die draaien op IBM Z-hardware, vaak ingezet voor moderne middleware, microservices en integratie met open-source ecosystemen \autocite{IBMRedbook_zOSBasics}.
	
	\item \textbf{z/VSE:} Een eenvoudiger mainframebesturingssysteem dat vooral gebruikt wordt in kleinere omgevingen, met ondersteuning voor batch- en transactionele workloads \autocite{IBMRedbook_zOSBasics}.
	
	\item \textbf{KVM:} Een open-source virtualisatielaag voor Linux workloads op IBM Z en LinuxONE, als alternatief voor z/VM in specifieke scenario’s \autocite{ibm_kvm_ibmz_linuxone}.
\end{itemize}

\section{\IfLanguageName{dutch}{Welke sectoren gebruiken mainframes?}{Which sectors use mainframes?}}
\label{sec:welke-sectoren-gebruiken-mainframes}

Mainframes worden voornamelijk ingezet in sectoren waar betrouwbaarheid, schaalbaarheid en beveiliging van cruciaal belang zijn. Mainframes worden wereldwijd gebruikt in onder meer de volgende sectoren:

\begin{itemize}

\item \textbf{Financiële sector:} Mainframes zijn de technologische ruggengraat van de globale financiële sector. 44 van de top 50 banken wereldwijd implementeren IBM mainframes voor de verwerking van miljarden dagelijkse transacties~\autocite{ibm_case_bankwest}. 
\item \textbf{Overheidsinstanties:} Overheden maken gebruik van mainframes voor belastingverwerking, sociale zekerheid, burgeradministratie en andere kritieke publieke diensten. Daarnaast vertrouwen overheidsorganisaties, waaronder rechtshandhavingsinstanties en nationale veiligheidsdiensten, op mainframes vanwege hun combinatie van beveiliging, prestaties en veerkracht. In deze context is bescherming tegen systeemuitval en beveiligingsincidenten van bijzonder belang. Moderne mainframeplatformen integreren bovendien geavanceerde analysemogelijkheden en AI-functionaliteiten om gegevens efficiënter te verwerken en cybersecurityprocessen te versterken~\autocite{IBMMainframe}.
\item \textbf{Gezondheidszorg:} Ziekenhuizen, zorgverzekeraars en andere gezondheidsorganisaties gebruiken mainframes voor het verwerken en beheren van grote hoeveelheden gevoelige gegevens, waaronder persoonlijk identificeerbare informatie (PII), medische dossiers en facturatiegegevens. Vooral grote internationale verzekeringsmaatschappijen vertrouwen op mainframes om deze data veilig en betrouwbaar te verwerken. In deze sector zijn gegevensintegriteit, privacy en beveiliging van essentieel belang~\autocite{IBMMainframe}.
\item \textbf{Detailhandel:} Grote retailorganisaties vertrouwen op mainframes voor voorraadbeheer, point-of-sale transactieverwerking en analyse van klantgegevens, vooral tijdens piekperiodes zoals promoties en feestdagen. Ook online retailers maken gebruik van mainframes om grote hoeveelheden transacties via mobiele en digitale kanalen op schaal te verwerken. De verwerkingscapaciteit van mainframes ondersteunt e-commerceplatformen die continu beschikbaar moeten zijn en hoge transactiesnelheden vereisen~\autocite{IBMMainframe}.
\item \textbf{Luchtvaartmaatschappijen:} Mainframes ondersteunen reserveringssystemen, vluchtplanning en passagiersbeheer. De mogelijkheid om grote aantallen transacties gelijktijdig te verwerken is hierbij essentieel ~\autocite{IBMMainframe}.
\item \textbf{Telecommunicatie:} Telecombedrijven gebruiken mainframes voor het beheren van klantendatabases, facturatie en netwerkbeheer, waarbij hoge beschikbaarheid en real-time verwerking noodzakelijk zijn ~\autocite{IBMMainframe}.
\item \textbf{Overige sectoren:} Ook in sectoren zoals energie, transport en productie (bijvoorbeeld, Arcelor Mittal) worden mainframes ingezet voor grootschalige dataverwerking en kritieke operationele systemen.

\end{itemize}

\subsection{\IfLanguageName{dutch}{Kenmerken van mainframes}{Characteristics of mainframes}}
\label{sec:kenmerken-van-mainframe}
Alhoewel er alternatieven voor mainframes bestaan, zoals Linux- en Windows-servers, onderscheiden mainframes zich door hun unieke kenmerken die hen bijzonder geschikt maken voor bedrijfskritische toepassingen.

Uit het onderzoek van \textcite{candido2024_mainframe_cloud} blijven mainframes relevant en vaak zelfs voordeliger voor bedrijfskritische, transactionele workloads, terwijl cloud computing superieur is voor elastische, schaalbare en moderne applicaties. Beide hebben hun plaats binnen hedendaagse IT-infrastructuren.

\subsection{\IfLanguageName{dutch}{Beschikbaarheid en uptime}{Availability and uptime of mainframes}}
\label{sec:beschikbaarheid-en-uptime}
Mainframes zijn ontworpen om 24/7 te draaien zonder noemenswaardige onderbrekingen. De IBM Z-mainframefamilie heeft een uptime van 99,999997\%, wat betekent dat bedrijven die afhankelijk zijn van mainframes zelden tot nooit te maken hebben metsysteemuitval \autocite{IBMMainframe}.

Door het gebruik van concepten en technologieën zoals Parallel Sysplex kan deze beschikbaarheid verder worden verhoogd. Hierbij werken meerdere systemen samen als één logisch geheel, waardoor workloads automatisch kunnen worden overgenomen bij uitval van een component. Dit beperkt het risico op onderbrekingen tot een minimum. (zie Sectie 2.3.4)

Dit niveau van betrouwbaarheid is essentieel voor bedrijfskritieke operaties, zoals het verwerken van financiële transacties of het beheren van luchtvaartreserveringen.

\subsection{\IfLanguageName{dutch}{Verwerkingskracht}{Processing power of mainframes}}
\label{sec:verwerkingskracht}
Mainframes zijn uitgerust met krachtige CPU-architecturen, zoals de S/390x, die meer dan 13 keer sneller zijn dan gemiddelde Intel-processors \autocite{IBMMainframe}. 

Dit stelt mainframes in staat om immense workloads aan te kunnen zonder prestatie verlies. 

Bovendien kunnen mainframes honderden virtuele machines parallellaten draaien, wat zorgt voor een flexibele en efficiënte inzet van IT-middelen.

\subsection{\IfLanguageName{dutch}{AI-capaciteit}{AI capacity of mainframes}}
\label{sec:ai-capaciteit}
Moderne mainframes zijn uitgerust met ingebouwde AI-chips, waardoor AI-modellen dicht bij de data kunnen worden verwerkt, zonder dat gegevens eerst moeten worden verplaatst naar een externe AI-server \autocite{IBMMainframe}. 

Dit maakt mainframes ideaal voor het uitvoeren van complexe AI-workloads, zoals risicoanalyse en fraudedetectie.

\subsection{\IfLanguageName{dutch}{Beveiliging}{Security and military context}}
\label{sec:beveiliging-militaire-context}
Mainframes zijn uitgerust met geavanceerde beveiligingsfuncties, zoals quantum-safe encryptie, die ervoor zorgen dat gegevens veilig blijven,zelfs in het tijdperk van quantum computing \autocite{IBMMainframe}. 

Bovendien staan compliance en beveiliging al vijf jaar op rij bovenaan de lijst van prioriteiten voor mainframe-gebruikers \autocite{BMC2024}.

In omgevingen waar nationale veiligheid en defensie centraal staan, zijn vertrouwelijkheid, integriteit en beschikbaarheid van gegevens van strategisch belang. Overheids- en defensieorganisaties verwerken grote hoeveelheden gevoelige of geclassificeerde informatie, waarbij strikte toegangscontrole en traceerbaarheid essentieel zijn. Mainframes worden in dergelijke contexten ingezet vanwege hun sterke isolatie tussen workloads, gecentraliseerde beveiligingsarchitectuur en ondersteuning voor geavanceerde encryptietechnologieën. \autocite{IBM_zOS_Security}

\section{\IfLanguageName{dutch}{Eigenschappen van z/OS architectuur}{Characteristics of z/OS architecture}}
\label{sec:kenmerken-van-zos-architectuur}
Deze sectie beschrijft sommige z/OS-concepten die nodig zijn om z/OS omgeving en taken van SysAdmin/SysProg te begrijpen.
\subsection{\IfLanguageName{dutch}{Datasets}{Datasets on z/OS}}
\label{sec:datasets-op-zos}
Binnen z/OS wordt data beheerd via zogenaamde data sets, wat de mainframe equivalent is van bestanden. Een data set bestaat uit records (vaste of variabele records) die door programma’s worden verwerkt. Datasets worden doorgaans aangesproken via hun datasetnaam (DSN); wanneer ze cataloged zijn, kan het systeem hun fysieke locatie automatisch terugvinden via de catalogus.
z/OS ondersteunt meerdere datasettypes naargelang de access method en het gebruik:
\begin{itemize}
	\item \textbf{Sequential (PS):} records worden sequentieel opgeslagen en gelezen; typisch voor batchverwerking.
	\item \textbf{Partitioned (PDS/PDSE):} bibliotheekstructuur met members; vaak gebruikt voor programmabibliotheken en procedurelibraries.
	\item \textbf{VSAM:} geoptimaliseerd voor efficiënte toegang; o.a. KSDS (key-based), ESDS (sequentieel), RRDS (relatief recordnummer) en LDS (byte-stream).
\end{itemize}

Data sets worden doorgaans opgeslagen op DASD-schijven of tape en worden beheerd via specifieke access methods die bepalen hoe data fysiek wordt opgeslagen en opgehaald. \autocite{IBM_zOS_Concepts}

\subsection{\IfLanguageName{dutch}{EBCDIC tekencodering}{EBCDIC on z/OS}}
\label{sec:ebcdic-op-zos}

EBCDIC (Extended Binary Coded Decimal Interchange Code) is de traditionele tekencodering binnen z/OS-datasets. In tegenstelling tot de wereldwijd dominante ASCII-standaard gebruiken klassieke mainframeomgevingen EBCDIC voor de representatie van karakters. Beide zijn 8-bit tekensets, maar verschillen in binaire toewijzing en sorteervolgorde, wat impact heeft op tekstverwerking en gegevensuitwisseling. Bij communicatie tussen ASCII- en EBCDIC-gebaseerde systemen is correcte conversie essentieel; bij binaire overdracht moet expliciet een binaire transfer worden gekozen om ongewenste karakterconversie te vermijden. \autocite{IBM_zOS_Concepts}

\subsection{\IfLanguageName{dutch}{Address spaces en virtual storage}{Address spaces}}
\label{sec:adress-spaces}

z/OS maakt gebruik van virtual storage, waarbij enkel actieve delen van een address space in het centrale geheugen worden geladen en inactieve delen tijdelijk op schijf worden geplaatst. Dankzij 64-bit adressering kan het systeem een zeer grote hoeveelheid virtueel geheugen aanspreken, aanzienlijk meer dan fysiek aanwezig is. \autocite{IBM_zOS_Concepts}

Een address space in z/OS is het bereik van virtuele adressen dat aan een gebruiker, job of subsystem wordt toegewezen. Elke address space is logisch geïsoleerd, waardoor programma’s en data van elkaar gescheiden blijven. Binnen één address space kunnen meerdere taken gelijktijdig worden uitgevoerd, wat multiprogramming mogelijk maakt. Fouten blijven doorgaans beperkt tot de betrokken address space, wat de stabiliteit en beveiliging van het systeem verhoogt. \autocite{IBM_zOS_Concepts}


\subsection{\IfLanguageName{dutch}{Parallel Sysplex}{Parallel Sysplex}}
\label{sec:parallel-sysplex}

IBM Parallel Sysplex is een clusteringtechnologie waarmee meerdere z/OS-systemen samenwerken als één logisch platform. Hierdoor kunnen workloads en resources gedeeld en dynamisch verdeeld worden, wat de schaalbaarheid en beschikbaarheid aanzienlijk verhoogt. \autocite{ibm_parallel_sysplex}

\subsection{\IfLanguageName{dutch}{Coupling Facility}{Coupling Facility}}
\label{sec:coupling-facility}

De Coupling Facility (CF) is een kerncomponent binnen Parallel Sysplex die gedeelde data en statusinformatie tussen meerdere z/OS-systemen beheert. Ze maakt coherente en veilige gegevensuitwisseling mogelijk zonder inconsistentie of corruptie, wat essentieel is voor hoge beschikbaarheid en schaalbaarheid. \autocite{ibm_parallel_sysplex}

De CF kan als aparte LPAR of als externe hardware worden ingezet en biedt diensten zoals locking, caching en list-processing. Hierdoor kunnen subsystemen zoals databases en message queues efficiënt samenwerken over meerdere systemen heen.

\subsection{\IfLanguageName{dutch}{Workload Manager (WLM)}{Workload Manager (WLM)}}
\label{sec:workload-management}

De Workload Manager (WLM) is een kerncomponent van z/OS die systeemresources dynamisch verdeelt over verschillende workloads. In goal mode gebeurt deze verdeling automatisch op basis van vooraf gedefinieerde servicedoelstellingen, zoals responstijd en prioriteit. \autocite{ibm_zos_wlm}

WLM ondersteunt zowel traditionele MVS-workloads (zoals batch en TSO) als z/OS UNIX-processen en draagt bij aan optimale performantie, resourcebenutting en stabiliteit binnen één systeem of over meerdere systemen in een sysplex. \autocite{ibm_zos_wlm}

\subsection{\IfLanguageName{dutch}{TSO op z/OS}{TSO on z/OS}}
\label{sec:tso-op-zos}

TSO/E (Time Sharing Option/Extensions) biedt gebruikers een interactieve toegang tot z/OS via een commandogebaseerde sessie, in tegenstelling tot batchverwerking. Authenticatie gebeurt via RACF en elke gebruiker werkt binnen een eigen sessie. \autocite{IBM_zOS_Concepts}

In de praktijk wordt TSO meestal gebruikt in combinatie met ISPF, een menu- en paneelgebaseerde interface die het beheer van datasets en programmabewerking vereenvoudigt. Toegang gebeurt via 3270-terminals of TN3270-emulators. \autocite{IBM_zOS_Concepts}

\subsection{\IfLanguageName{dutch}{ISPF (Interactive System Productivity Facility)}{ISPF (Interactive System Productivity Facility)}}
\label{sec:ispf}

ISPF (Interactive System Productivity Facility) is een menu- en paneelgebaseerde werkomgeving bovenop TSO/E en vormt voor veel gebruikers de primaire interface binnen z/OS. Het biedt functionaliteiten zoals datasetbeheer, tekstbewerking en toegang tot diverse systeemutilities, waardoor interactief systeembeheer efficiënter verloopt. \autocite{IBM_zOS_Concepts}

ISPF is hiërarchisch opgebouwd via een Primary Option Menu en kan per installatie worden aangepast. Hoewel traditioneel toetsenbordgestuurd via 3270-emulatie, bestaan er ook moderne GUI-clients. \autocite{IBM_zOS_Concepts}

\subsection{\IfLanguageName{dutch}{z/OS UNIX System Services}{z/OS UNIX System Services}}
\label{sec:z-os-unix-system-services}

z/OS UNIX System Services (USS) biedt een POSIX-conforme UNIX-omgeving binnen z/OS. Het stelt gebruikers in staat om via een shell UNIX-commando’s uit te voeren, scripts te ontwikkelen en applicaties te starten, waardoor integratie met open systemen en moderne ontwikkelomgevingen mogelijk wordt. \autocite{IBM_zOS_Concepts}

USS kan worden benaderd via 3270-emulatie, TCP/IP of vanuit TSO (OMVS) en vormt een brug tussen traditionele mainframeworkloads en UNIX-gebaseerde toepassingen.

\subsection{\IfLanguageName{dutch}{Beveiliging in z/OS (RACF)}{Security in z/OS}}
\label{sec:network-and-security}

RACF (Resource Access Control Facility) is de centrale beveiligingscomponent van z/OS en beheert authenticatie en autorisatie van gebruikers en systeemresources. Op basis van beveiligingsprofielen bepaalt RACF welke gebruikers toegang krijgen tot datasets, applicaties en systeemfuncties. \autocite{IBM_zOS_Concepts}

Het ondersteunt het principe van least privilege, logging van toegangsacties en beleidsregels rond wachtwoorden en accountbeheer. Subsystemen zoals CICS en Db2 integreren met RACF voor consistente toegangscontrole binnen de mainframeomgeving. \autocite{IBM_zOS_Concepts}

\subsection{\IfLanguageName{dutch}{Job Control Language (JCL)}{Job Control Language (JCL)}}
\label{sec:jcl}

Job Control Language (JCL) is de taal waarmee gebruikers aan z/OS specificeren hoe een batchjob moet worden uitgevoerd. Via JCL wordt bepaald welk programma moet draaien, welke datasets als input of output worden gebruikt en hoe de uitvoer verwerkt moet worden. \autocite{IBM_zOS_Concepts}

Een JCL-job bestaat hoofdzakelijk uit drie basisstatements: JOB, EXEC en DD. Het JOB-statement definieert de workload, EXEC start een programma of procedure, en DD (Data Definition) koppelt datasets of I/O-resources aan het programma. \autocite{IBM_zOS_Concepts}

JCL wordt voornamelijk gebruikt voor batchverwerking en vormt een essentieel onderdeel van operationeel systeembeheer binnen z/OS-omgevingen.

\subsection{\IfLanguageName{dutch}{REXX}{REXX on z/OS}}
\label{sec:rexx-op-zos}

REXX (Restructured Extended Executor) is een procedurele programmeertaal die op z/OS voornamelijk wordt gebruikt voor automatisering en scripting. De taal staat bekend om haar leesbaarheid en eenvoudige syntaxis en wordt meestal geïnterpreteerd uitgevoerd, hoewel compilatie mogelijk is om de performantie te verbeteren \autocite{IBM_zOS_Concepts}.

Binnen z/OS wordt REXX vaak ingezet voor het uitvoeren van TSO/E-commando’s, het automatiseren van beheertaken, het aanroepen van andere programma’s en het ontwikkelen van ISPF-applicaties. Hierdoor vormt REXX een belangrijk hulpmiddel voor systeembeheerders bij operationele taken en procesautomatisering.

\subsection{\IfLanguageName{dutch}{SMP/E}{SMP/E on z/OS}}
\label{sec:smpe-op-zos}
System Modification Program Extended (SMP/E) is de standaardtool binnen z/OS voor het installeren, onderhouden en beheren van softwareproducten. Daarnaast houdt SMP/E alle wijzigingen en updates aan deze software systematisch bij \autocite{IBM_zOS_Concepts}.

\subsection{\IfLanguageName{dutch}{Subsystemen in z/OS}{Subsystems in z/OS}}
\label{sec:fsubsystemen-in-zos}

Deze zijn gespecialiseerde subsystemen in z/OS en een SysAdmin/SysProg moet deze subsystemen weten en begrijpen. In deze bachelorproef gaat er minder aandacht besteed worden op deze subsystemen.
\begin{itemize}
	\item \textbf{CICS:} Customer Information Control System voor transactionele toepassingen
	\item \textbf{IMS DB:} Information Management System voor databasebeheer
	\item \textbf{IMS TM:} Information Management System voor transactiebeheer
	\item \textbf{DB2:} relationele database
	\item \textbf{MQ:} message queuing voor communicatie tussen applicaties
	\item \textbf{IWS}: Integrated Workload Scheduler voor geavanceerde jobplanning en -beheer
\end{itemize}
Deze subsystemen draaien binnen z/OS maar beheren elk hun eigen type workload (batch, transacties, databasebeheer).



\section{\IfLanguageName{dutch}{Interactie met z/OS}{Interacting with z/OS}}
\label{sec:interactie-met-zos}
z/OS biedt verschillende manieren waarop gebruikers en systeembeheerders met het systeem kunnen interageren.

\paragraph{Traditionele toegang}
Historisch gezien gebeurde interactie met z/OS via fysieke 3270-terminals 
die rechtstreeks verbonden waren met het mainframe. Gebruikers bevonden zich 
in hetzelfde datacenter of maakten gebruik van dedicated netwerkverbindingen 
naar het systeem. De terminal fungeerde als invoer en uitvoerapparaat, 
waarmee sessies in TSO, ISPF of CICS konden worden gestart.

Deze vorm van interactie vereiste fysieke infrastructuur en directe 
netwerkconnectiviteit met het mainframe. Hoewel fysieke terminals vandaag 
grotendeels vervangen zijn door softwarematige emulators.

\paragraph{Emulators en externe toegang}
Fysieke 3270-terminals zijn vervangen door softwarematige 3270-emulators die via het TN3270-protocol 
verbinding maken met het mainframe. Daarnaast ondersteunt z/OS via UNIX System Services (USS) 
toegang via SSH. Voor developers wordt vaak gebruikgemaakt van IBM Developer for z/OS (IDz), 
een Eclipse gebaseerde ontwikkelomgeving die integratie biedt met z/OS. 

\paragraph{API’s en moderne integratie}
Moderne z/OS-omgevingen ondersteunen steeds vaker API-gebaseerde interactie. 
Via z/OS Management Facility (z/OSMF) worden RESTful services aangeboden 
die programmatische toegang tot datasets, jobs en systeemfuncties mogelijk maken. 
Dit maakt integratie met automatiseringstools en DevOps-pipelines mogelijk.

\subsection{Modernisering en hybride integratie van z/OS}
\label{sec:modernisering-integratie-zos}

De modernisering van z/OS-omgevingen richt zich niet op het vervangen van het mainframe, maar op het integreren ervan binnen hybride en cloudgeoriënteerde IT-architecturen. Organisaties streven naar een model waarbij bestaande, bedrijfskritieke toepassingen behouden blijven, terwijl nieuwe functionaliteiten worden ontwikkeld met behulp van moderne tools en standaarden \autocite{ibm2023hybrid}.

Een belangrijke rol in deze evolutie wordt gespeeld door z/OS Management Facility (z/OSMF), dat een webgebaseerde beheerinterface biedt voor configuratie, monitoring en workflowautomatisering binnen z/OS \autocite{ibm2023zosmf}. z/OSMF ondersteunt self-service en REST-gebaseerde interacties, waardoor traditionele beheertaken beter aansluiten bij hedendaagse IT-praktijken.

Daarnaast vormt Zowe een open-source initiatief dat z/OS toegankelijk maakt via gestandaardiseerde API’s, command line interfaces en integraties met moderne ontwikkelomgevingen \autocite{zowe2023}. Door deze abstrahering wordt de kloof tussen klassieke mainframe-omgevingen en cloud-native ontwikkelmodellen verkleind.

De combinatie van dergelijke technologieën maakt het mogelijk om z/OS-systemen op een gecontroleerde en veilige manier te integreren in hybride cloudarchitecturen, waarbij stabiliteit en innovatie hand in hand gaan.


\section{\IfLanguageName{dutch}{Belang van z/OS in mainframes}{Role of z/OS in mainframes}}
\label{sec:role-van-zos-in-mainframes}
z/OS vormt het centrale besturingssysteem van IBM Z-mainframes en fungeert als het operationele hart van de mainframeomgeving. Het beheert de hardwarebronnen zoals CPU, geheugen, opslag (DASD) en I/O-apparaten, en coördineert de gelijktijdige uitvoering van duizenden programma’s en gebruikerssessies.\autocite{IBM_zOS_Concepts}

\subsection{Wat is system administration?}
\label{sec:wat-is-systembeheer}

Systeembeheer (of system administration) verwijst naar het beheren, configureren, monitoren en onderhouden van computersystemen en IT-infrastructuur binnen een organisatie. Het doel van systeembeheer is het waarborgen van beschikbaarheid, veiligheid, prestaties en continuïteit van systemen en diensten. \autocite{TechNeeds2025}

Een systeembeheeromgeving bestaat doorgaans uit drie kernaspecten:
\begin{itemize}
	\item \textbf{Systeemconfiguratie en onderhoud:} Installeren, updaten en configureren van besturingssystemen, software en hardwarecomponenten.
	\item \textbf{Monitoring en probleemoplossing:} Toezicht houden op systeemprestaties, detecteren van incidenten en oplossen van technische problemen.
	\item \textbf{Beveiliging en toegangsbeheer:} Beheren van gebruikersaccounts, toegangsrechten en beveiligingsmaatregelen om systemen te beschermen tegen ongeautoriseerde toegang.
\end{itemize}

Zonder effectief systeembeheer lopen organisaties het risico op systeemuitval, beveiligingsincidenten en prestatieproblemen, wat directe impact kan hebben op bedrijfsprocessen en dienstverlening.

\subsection{Waarom is system administration belangrijk?}
\label{sec:waarom-sysbeheer-belang}

Gezien het z/OS besturingsysteem letterlijk het hart is van de mainframes en bedrijfskritische subsystemen zoals Db2, IMS en CICS draaien op z/OS is de correcte werking van deze subsystemen rechtstreeks afhankelijk van de stabiliteit, configuratie en beveiliging van het onderliggende z/OS-platform.

System Administration speelt hierin een cruciale rol. De system administrator staat in voor het onderhoud van systeemsoftware, het toepassen van updates en patches, het monitoren van prestaties en het bewaken van beschikbaarheid. Indien software-updates niet tijdig worden uitgevoerd of configuraties foutief worden aangepast, kan dit leiden tot beveiligingsrisico’s, prestatieproblemen of zelfs productie-uitval. \autocite{IBM_zOS_Concepts}

Daarnaast is beveiliging een fundamenteel onderdeel van systeembeheer. Via tools zoals RACF, auditing en change management wordt gewaarborgd dat enkel geautoriseerde gebruikers toegang hebben tot gevoelige data en systeemresources. In sectoren zoals financiën, overheid en gezondheidszorg kan falend systeembeheer ernstige operationele en juridische gevolgen hebben. \autocite{IBM_zOS_Concepts}

Effectief systeembeheer is bijgevolg essentieel om continuïteit, betrouwbaarheid en compliance binnen een mainframeomgeving te garanderen.

\subsection{\IfLanguageName{dutch}{Wat is het verschil tussen system administrator en system programmer}{What is the role of a system administrator in a z/OS context}}
\label{sec:verschil-sysadmin-sysprog}
De system programmer beheert de installatie, configuratie en technische optimalisatie van z/OS en systeemsoftware, inclusief maintenance, tuning en complexe probleemanalyse. \autocite{IBM_zOS_Concepts}
De system administrator beheert de dagelijkse werking van de z/OS-omgeving, inclusief toegangsbeheer, monitoring en operationele ondersteuning van applicaties en data.\autocite{IBM_zOS_Concepts}
Maar het onderscheid varieert tussen een system programmer en een system administrator afhankelijk van de grootte en structuur van de organisatie. In kleinere IT-omgevingen kunnen beide rollen door één persoon worden uitgevoerd. In grotere organisaties worden de verantwoordelijkheden doorgaans gescheiden omwille van efficiëntie en auditvereisten. 
Hoewel de focus van de system administrator operationeel is, vertoont de rol in de praktijk gedeeltelijke overlap met die van de system programmer. Relevante systeemgerichte taken worden daarom meegenomen in dit competentieframework zoals het werken met SMP/E software beheer tool.

\subsection{\IfLanguageName{dutch}{Wat zijn de taken van een system administrator volgens IBM}{What is the role of a system administrator in a z/OS context}}
\label{sec:rol-van-een-system-administrator-in-zos-context}
De system administrator is in hoofdzaak verantwoordelijk voor het dagelijkse operationele beheer van de mainframeomgeving vanuit een bedrijfsgericht perspectief. Waar de system programmer zich richt op het onderhoud van het besturingssysteem en de onderliggende infrastructuur, concentreert de system administrator zich op het beheer van bedrijfskritische data, applicaties en toegangscontrole. \autocite{IBM_zOS_Concepts}

Volgens \autocite{IBM_zOS_Concepts} typische taken van een system administrator omvatten:

\begin{itemize}
	\item Installeren en configureren van softwarecomponenten
	\item Beheren van gebruikersaccounts en gebruikersprofielen
	\item Onderhouden van beveiligings- en toegangsrechten
	\item Beheren van opslagmiddelen, printers en randapparatuur
	\item Beheren van netwerkconnectiviteit
	\item Monitoren van systeemprestaties en beschikbaarheid
\end{itemize}

Daarnaast werkt de system administrator nauw samen met applicatieontwikkelaars en eindgebruikers om te verzekeren dat toepassingen correct functioneren binnen de operationele omgeving. In grotere organisaties kunnen gespecialiseerde rollen bestaan, zoals database administrators (DBA’s) of security administrators, die specifieke onderdelen van het beheer uitvoeren. \autocite{IBM_zOS_Concepts}

\subsection{De evoluerende rol van de SysAdmin/SysProg}
\label{sec:evolution-sys-admin}

De rol van de System Administrator heeft in de loop der jaren een duidelijke evolutie doorgemaakt. Waar system administrators traditioneel voornamelijk verantwoordelijk waren voor het installeren, configureren en onderhouden van servers en besturingssystemen, omvat hun takenpakket vandaag een bredere set aan verantwoordelijkheden die automatisering, beveiliging en cloudintegratie omvatten \autocite{kim2016devops}.

De opkomst van cloud computing heeft deze verschuiving verder versterkt. Volgens NIST wordt cloud computing gedefinieerd als een model dat on-demand toegang biedt tot gedeelde en configureerbare IT-resources \autocite{mell2011nist}. Deze evolutie heeft geleid tot een overgang van handmatig infrastructuurbeheer naar geautomatiseerde en schaalbare beheermodellen.

Daarnaast tonen recente industriële rapporten aan dat organisaties steeds meer inzetten op containerisatie, hybride cloudarchitecturen en DevOps-praktijken \autocite{redhat2023state}. Hierdoor wordt van System Administrators verwacht dat zij niet enkel technische expertise bezitten, maar ook competenties ontwikkelen in automatisering, samenwerking en continue verbetering.

Binnen mainframe- en z/OS-omgevingen vertaalt deze evolutie zich in een combinatie van traditionele stabiliteit en moderne integratie, waarbij system administrators een sleutelrol spelen in het veilig en efficiënt verbinden van bestaande infrastructuur met hedendaagse IT-architecturen.

\subsection{Focus op z/OS, JCL, REXX en RACF in deze bachelorproef}
\label{sec:focus-op-jcl-rexx}

Gezien de blijvende relevantie van mainframe-omgevingen binnen kritieke sectoren zoals financiën en overheid, richt deze bachelorproef zich op z/OS en de kerncomponenten JCL, REXX en RACF. De redenen hiervoor zijn:

\begin{itemize}
	\item \textbf{Blijvende relevantie van mainframes:} z/OS blijft een essentieel platform voor mission-critical toepassingen waar stabiliteit, betrouwbaarheid en veiligheid centraal staan.
	
	\item \textbf{Operationele basiskennis:} JCL en REXX vormen de fundamenten van dagelijkse beheer- en automatisatietaken binnen z/OS.
	
	\item \textbf{Beveiliging en governance:} RACF speelt een cruciale rol in toegangsbeheer, compliance en risicobeperking binnen enterprise-omgevingen.
	\item \textbf{Ondersteuning van modernisering:} Een solide kennis van deze kerncomponenten is noodzakelijk om mainframe-omgevingen succesvol te integreren in hybride en gemoderniseerde IT-architecturen.
	\item \textbf{Software onderhoud en installatie:} Aangezien veel bedrijven gebruik maken van verschillende software tools van allerlei vendors is het belangrijk om SMP/E te begrijpen en dit te kunnen gebruiken voor een Level 2 System Administrator.
\end{itemize}

\section{Uitdagingen in de mainframe-sector: vergrijzing en kennisoverdracht}
\label{sec:mainframe-uitdagingen-vergrijzing}

Ondanks de blijvende relevantie van mainframes kampen organisaties met een structurele personeelsuitdaging. De instroom van nieuwe profielen blijft beperkt, onder meer omdat mainframes minder zichtbaar is binnen klassieke IT-opleidingen en veel starters eerder richting cloud-native stacks en populaire programmeertalen worden getrokken. Hierdoor ontstaat een krapte op de arbeidsmarkt, zeker voor profielen die zowel de z/OS-basis als enterprise-processen beheersen.

\subsubsection{Conclusie rond de uitdagingen}
\label{sec:uitdagingen-standvanzaken-conclusie}
De literatuur toont duidelijk aan dat er een structureel tekort is aan mainframe-expertise en dat bestaande opleidingsinitiatieven onvoldoende schaalbaar zijn.
Daarom neemt dit onderzoek deze problematiek als uitgangspunt en vertaalt deze naar concrete competentievereisten binnen het voorgestelde framework.
Daarnaast wordt de focus gelegd op praktijkgerichte en meetbare vaardigheden, om beter aan te sluiten bij de noden van de industrie. Wat na gesprekken en interviews met de professionals meer verfijnd zullen worden.

Een belangrijke uitdaging is de veroudering van het bestaande personeelsbestand. Volgens \textcite{astadia_retirement} is ongeveer 60\% van de mainframeprofessionals ouder dan 50 jaar, wat impliceert dat een aanzienlijk deel van deze expertise op middellange termijn de arbeidsmarkt zal verlaten.

Deze demografische evolutie leidt tot een groeiend tekort aan gekwalificeerde professionals om openstaande functies in te vullen. Het probleem wordt versterkt door de beperkte instroom van nieuwe talenten. Uit een studie blijkt dat 64\% van de IT-operationsmanagers aangeeft dat het vinden van mainframespecialisten een significante uitdaging vormt. \autocite{astadia_retirement}

\subsection{Impact op organisaties en mogelijke oplossingen}
\label{sec:bestaande-ops}

De toenemende vergrijzing en het tekort aan gespecialiseerde mainframeprofessionals hebben geleid tot verschillende initiatieven vanuit de industrie. Organisaties investeren in gerichte opleidingsprogramma’s, reskilling-trajecten en samenwerking met onderwijsinstellingen om nieuwe instroom te stimuleren. Volgens Kyndryl \autocite{kyndrylskillsgap} vereist het dichten van de mainframe skills gap een combinatie van formele training, mentoring en modernisering van het imago van mainframe-technologie binnen bredere IT-contexten. 

\subsection{De rol van onderwijs in mainframe opleiding}
\label{sec:mainframe-opleidingen}

Onderwijsinstellingen spelen een cruciale rol in het aanpakken van de instroomproblematiek. Door mainframeconcepten te integreren in bestaande IT-curricula kunnen studenten vroegtijdig kennismaken met mainframes. Industry academia samenwerking, gastcolleges en praktijkgerichte labs worden in dit verband vaak aangehaald als effectieve strategieën \autocite{kyndrylskillsgap}.

\subsection{HOGENT’s unieke aanbod}

HOGENT onderscheidt zich door als enige instelling in de Benelux een specifieke opleiding tot mainframe-expert aan te bieden. Deze micro-credential richt zich tot IT-professionals en afgestudeerden met een bachelor- of masterdiploma in een IT-gerelateerde opleiding die hun expertise op korte termijn willen verdiepen en uitbreiden \autocite{hogent2025}.


\subsection{De kloof tussen onderwijs en industrie}

Ondanks het bestaan van verschillende competentieraamwerken blijft er een aanzienlijke kloof bestaan tussen wat onderwijsinstellingen aanbieden en wat de industrie nodig heeft, vooral op het gebied van mainframe-technologieën. Deze kloof wordt versterkt door verschillende factoren:

\begin{itemize}
	\item \textbf{Beperkte toegang tot mainframe-technologie:} Veel onderwijsinstellingen hebben geen directe toegang tot mainframe-systemen voor praktijkonderwijs, waardoor het moeilijk is om hands-on ervaring aan te bieden.
	
	\item \textbf{Gebrek aan gespecialiseerde docenten:} Er is een tekort aan docenten met actuele mainframe-expertise, mede door de hoge vraag naar deze professionals in de industrie.
	
	\item \textbf{Snelle technologische evolutie:} Mainframe-technologieën blijven evolueren, wat het voor onderwijsprogramma’s uitdagend maakt om curricula voortdurend up-to-date te houden.
	
	\item \textbf{Perceptieproblemen:} Er bestaat een hardnekkige perceptie dat mainframes verouderde technologie zijn, wat studenten kan ontmoedigen om zich in deze richting te specialiseren.
\end{itemize}

Het overbruggen van deze kloof vereist een gecoördineerde inspanning van zowel onderwijsinstellingen als de industrie. Initiatieven zoals het IBM Z Academic Initiative en het Open Mainframe Project bieden ondersteuning, leermaterialen en toegang tot platformen voor instellingen die mainframe-onderwijs willen integreren in hun curricula.
\subsection{Competency-Based Education (CBE) als opleidings- en evaluatiekader}
\label{sec:cbe}

Competency-Based Education (CBE) is een onderwijsbenadering waarbij leerdoelen expliciet worden gedefinieerd als competenties, en waarbij voortgang wordt bepaald door aantoonbare beheersing in plaats van door contacturen of tijdsgebonden curricula. CBE sluit hierdoor nauw aan bij beroepsgerichte IT-profielen, waar competenties best zichtbaar worden via performantie in realistische taken \autocite{mulder2014competence}. 

Binnen dit onderzoek fungeert CBE als brug tussen het competentieframework (Level 1 en Level 2) en de opleiding: competenties worden vertaald naar observeerbare leeruitkomsten, oefenactiviteiten en toetscriteria (rubrics). Evaluatie gebeurt bij voorkeur op basis van authentieke bewijsvormen (bijv.\ uitgevoerde JCL-taken, gedocumenteerde RACF-wijzigingen, incidentanalyses), in lijn met de evidence-gerichte opbouw van het IBM apprenticeship model \autocite{ibm_level1_framework,ibm_level2_framework}. 

De niveaudifferentiatie binnen CBE wordt ondersteund door internationale raamwerken zoals SFIA en e-CF, waarin autonomie, complexiteit en verantwoordelijkheid expliciet toenemen met het competentieniveau \autocite{sfia_framework,eCF_framework}.

\subsection{Naar een geïntegreerd competentieraamwerk}

Een effectief competentieraamwerk voor System Administrators (Level 1 en Level 2) binnen een z/OS-omgeving moet rekening houden met de unieke kenmerken van mainframe-technologie, de specifieke vereisten van de industrie en de pedagogische principes die effectief leren ondersteunen.

Het framework dat in deze bachelorproef wordt ontwikkeld, streeft ernaar deze elementen te integreren in een coherent en praktisch toepasbaar model voor een onderwijscontext. Door gebruik te maken van bestaande raamwerken zoals het IBM Z System Administrator Competency Framework, het e-CF en SFIA, en deze te combineren met inzichten uit literatuur en praktijk, beoogt dit model een brug te slaan tussen theorie en industrie.

Daarnaast besteedt het framework expliciet aandacht aan competentieprogressie: van een beginnend, operationeel profiel (Level 1) naar een meer gevorderd profiel (Level 2). Dit sluit aan bij de niveaustructuren van SFIA en e-CF, waarin autonomie, complexiteit en verantwoordelijkheid gradueel toenemen \autocite{sfia_framework,eCF_framework}.

Door deze geïntegreerde benadering definieert het voorgestelde raamwerk niet alleen de technische competenties die nodig zijn voor effectief systeembeheer in een z/OS-omgeving, maar ook de professionele en gedragscompetenties die essentieel zijn voor duurzaam functioneren binnen moderne enterprise-IT-omgevingen.

\section{System Administrators en competentieontwikkeling}
\label{sec:sysadmin-comp-ontwkile}

System Administrators spelen een cruciale rol binnen z/OS door het waarborgen van stabiliteit, beveiliging en continuïteit van kritieke systemen. Organisaties zijn sterk afhankelijk van hun expertise voor het dagelijks beheer en de bescherming van deze infrastructuur.

Door de evolutie naar hybride cloudarchitecturen en verdere automatisering neemt de complexiteit van systeembeheer toe. Hierdoor worden hogere eisen gesteld aan zowel technische kennis als het vermogen om zich voortdurend aan te passen aan nieuwe technologieën. In deze sectie wordt onderzocht welke competenties noodzakelijk zijn en hoe gerichte competentieontwikkeling kan bijdragen aan toekomstbestendig systeembeheer.


\subsection{Modellen voor competentieontwikkeling: kennis, vaardigheden en professioneel gedrag}
\label{sec:ksa-modellen}

Competentieontwikkeling binnen ICT-profielen wordt doorgaans gestructureerd rond drie kerncomponenten: kennis (knowledge), vaardigheden (skills) en attitude of professioneel gedrag (attitude/behavior). Deze driedeling vormt een breed erkend model binnen zowel onderwijskundige literatuur als internationale competentieraamwerken \autocite{mulder2014competence}.

\paragraph{Kennis (Knowledge)}

Kennis verwijst naar het theoretisch en conceptueel begrip van systemen, architecturen en processen. Binnen het z/OS-domein omvat dit onder meer inzicht in mainframe-architectuur, I/O-structuren, adressering, subsystemen zoals CICS, IMS en Db2, en beveiligingsconcepten.

Zowel het IBM Level 1- als Level 2-competentieframework structureren expliciet kenniscomponenten als afzonderlijke beoordelingscriteria, bijvoorbeeld “Understanding of structured programming design”, “Understanding of z/OS system components” en “Understanding of security implementation” \autocite{ibm_level1_framework,ibm_level2_framework}.  

Ook binnen SFIA wordt kennis impliciet geïntegreerd in de beschrijving van vaardigheden per niveau, waarbij hogere niveaus een diepgaander conceptueel begrip vereisen \autocite{sfia_framework}. Het e-CF benoemt kennis expliciet als onderdeel van elke e-competentie, naast vaardigheden en autonomie \autocite{eCF_framework}.

\paragraph{Vaardigheden (Skills)}

Vaardigheden verwijzen naar het toepassen van kennis in concrete handelingen binnen een beroepscontext. Voor een z/OS System Administrator betreft dit onder meer:

\begin{itemize}
	\item Het coderen van JCL met correcte syntax
	\item Het ontwikkelen van REXX-scripts
	\item Het uitvoeren van SMP/E-onderhoud
	\item Het beheersen van RACF toegangsrechten
	\item Het begrijpen van z/OS architectuur
	\item Change management, incident analyse
\end{itemize}

De IBM-competentiedocumenten structureren deze vaardigheden expliciet als observeerbare prestaties met meetbare assessment criteria en evidence types (observatie, simulatie, mentorbeoordeling) \autocite{ibm_level1_framework,ibm_level2_framework}.  

Deze operationalisering sluit aan bij competentiegericht opleiden, waarbij aantoonbare taakuitvoering centraal staat \autocite{mulder2014competence}.

\paragraph{Attitude en professioneel gedrag (Behavior)}

Naast technische expertise benadrukken moderne competentieraamwerken het belang van professioneel gedrag. Dit omvat onder meer:

\begin{itemize}
	\item Methodisch en risicogestuurd werken
	\item Nauwkeurige documentatie
	\item Samenwerking en communicatie
	\item Feedbackcultuur en mentorschap
	\item Continue verbetering
\end{itemize}

Binnen het IBM-framework worden deze expliciet benoemd als gedragcompetenties (bijvoorbeeld teamwork, communication skills, feedback behavior en willingness to learn from mentors) \autocite{ibm_level1_framework}.  

\subsection{Integratie van educatieve theorieën in competentieontwikkeling voor z/OS System Administrators}
\label{sec:integratie-educatieve-theorie}

Het ontwikkelen van een competentieframework vereist niet alleen een beschrijving van vereiste kennis en vaardigheden, maar ook een onderbouwde visie op hoe deze competenties worden verworven en geëvalueerd \autocite{mulder2014competence}. Daarom wordt het framework didactisch verankerd in relevante educatieve theorieën.

\subsubsection{Volwassenenonderwijs en praktijkgericht leren}

Binnen professionele IT-opleidingen gaat het vaak om volwassen lerenden met voorkennis of werkervaring. Volgens principes van volwassenenonderwijs is leren effectiever wanneer het praktijkgericht, onmiddellijk toepasbaar en deels zelfgestuurd is. In de z/OS-context vertaalt dit zich naar realistische opdrachten, zoals het uitvoeren van JCL-jobs, analyseren van output via SDSF of het beheren van basis RACF-profielen met native RACF commando's.

\subsubsection{Ervaringsgericht leren en niveauprogressie}

Ervaringsgericht leren benadrukt het belang van leren via concrete uitvoering, reflectie en her-toepassing \autocite{kolb2015experiential}. Voor z/OS System Administrators impliceert dit een leertraject waarin technische handelingen in een labomgeving worden uitgevoerd, gevolgd door feedback en koppeling aan onderliggende systeemprincipes. 

Daarnaast speelt informele kennisoverdracht via mentorschap een belangrijke rol. In veel mainframe-omgevingen wordt expertise traditioneel doorgegeven via een senior–juniorrelatie, vergelijkbaar met het senpai–kōhai-principe, waarbij impliciete kennis, best practices en risicobeheersing in de praktijk worden overgedragen. Deze vorm van begeleide praktijkervaring ondersteunt de ontwikkeling van professionele oordeelsvorming en versnelt de overgang naar hogere competentieniveaus.

Deze aanpak ondersteunt de overgang van procedurele taakuitvoering (Level 1) naar analytische probleemoplossing en transfer naar nieuwe situaties (Level 2), wat aansluit bij de progressieniveaus zoals beschreven in SFIA \autocite{sfia_framework}.

\subsubsection{Cognitieve belasting en gefaseerde opbouw}

Gezien de complexiteit van z/OS-omgevingen is beheersing van cognitieve belasting essentieel \autocite{sweller2011cognitive}. Een gefaseerde opbouw van basisconcepten (TSO/ISPF, datasets, JCL) naar geïntegreerde scenario’s (subsystemen (CICS,IMS,Db2,MQ...), gevorderd REXX-automatisering, SMP/E-onderhoud) bevordert duurzame kennisverwerving en voorkomt overbelasting.

\subsubsection{Evaluatie en niveaudifferentiatie}

Competentiegericht opleiden vereist evaluatie in realistische beroepscontexten. In plaats van louter theoretische toetsing wordt gewerkt met observeerbare prestaties, zoals een correct uitgevoerde batchjob, een gedocumenteerde RACF-wijziging of een onderbouwde incidentanalyse. \autocite{sfia_framework,eCF_framework}

Het onderscheid tussen Level 1 en Level 2 wordt hierbij bepaald op basis van autonomie, complexiteit, risicobeheersing en kwaliteit van documentatie. Deze criteria sluiten aan bij internationale competentieraamwerken zoals SFIA en ondersteunen een objectieve positionering binnen het voorgestelde model.

\subsubsection{Didactische verankering van het framework}

Door de integratie van deze educatieve inzichten wordt het competentieframework meer dan een beschrijvend overzicht van technische vereisten. Het vormt een didactisch onderbouwd instrument dat inzetbaar is binnen zowel onderwijscontexten als professionele opleidingsprogramma’s.


\section{Vergelijking met andere competentieraamwerken}
\label{sec:vergelijking-andere-compframeworks}

Het opstellen van een competentieraamwerk is essentieel voor het definiëren van de vaardigheden, kennis en gedrag die nodig zijn voor specifieke functies binnen een organisatie. Dit helpt niet alleen bij het identificeren van lacunes in vaardigheden, maar ondersteunt ook bij werving, prestatiebeheer en loopbaanontwikkeling. 

In dit hoofdstuk vergelijken we verschillende bestaande competentieraamwerken die relevant zijn voor IT-professionals.
\subsection{Het IBM Z System Administrator Competency Framework – Level 1 en Level 2}

Het IBM Z System Administrator Competency Framework is een gestructureerd raamwerk ontwikkeld binnen IBM’s Apprenticeship Program als open standaard ter ondersteuning van opleidings- en werkplekleertrajecten voor mainframebeheerders \autocite{ibm_level1_framework,ibm_level2_framework}. 

Het raamwerk richt zich specifiek op de rol van IBM Z System Administrator en maakt een onderscheid tussen twee progressieniveaus: een operationeel instapniveau (Level 1) en een gevorderd, analytisch niveau (Level 2). Deze niveaus verschillen in mate van autonomie, complexiteit van taken en impact op productieomgevingen.

Het framework bestaat uit drie kerncomponenten:

\begin{itemize}
	\item \textbf{Werkprocessen:} Een beschrijving van de algemene vereisten die deelnemers moeten doorlopen, waaronder praktijkprojecten, mentoring, coaching, job shadowing en formele training.
	
	\item \textbf{Competenties en prestatiecriteria:} Een gedetailleerde opsplitsing van de vereiste technische kennis, vaardigheden en professioneel gedrag. Deze worden gestructureerd volgens kenniscompetenties (K), vaardigheidscompetenties (S) en gedragscompetenties (B).
	
	\item \textbf{Gerelateerde instructie:} Een overzicht van formele opleidingsmodules en leeractiviteiten die nodig zijn om de vereiste competenties te ontwikkelen en aan te tonen.
\end{itemize}

Binnen dit raamwerk worden kerncompetenties gedefinieerd die zowel technische expertise (zoals JCL, REXX, RACF en systeembeheer) als professionele vaardigheden (zoals communicatie, documentatie en risicobeheersing) omvatten. De overgang van Level 1 naar Level 2 wordt gekenmerkt door een verschuiving van begeleide, proceduregerichte uitvoering naar zelfstandige, analytische en risicogestuurde besluitvorming binnen productieomgevingen.

\subsubsection{Conclusie: Structuur van andere raamwerken copieeren}
De besproken modellen (kennis, vaardigheden en gedrag) tonen aan dat competentieontwikkeling meer omvat dan louter technische kennis.
Daarom neemt dit onderzoek de KSA-structuur (Knowledge, Skills, Attitude) over als basis voor het competentieframework,
en wordt deze verder geconcretiseerd naar specifieke z/OS-taken en observeerbare leeruitkomsten.


\subsection{Het Skills Framework for the Information Age (SFIA)}
\label{sec:sfia-framework}
Het Skills Framework for the Information Age (SFIA) is een wereldwijd erkend competentieraamwerk dat sinds 1999 in gebruik is. Het definieert 86 vaardigheden, elk beschreven over zeven niveaus van bekwaamheid. \autocite{sfia_framework}
SFIA is ontworpen om organisaties te helpen bij het identificeren van de vaardigheden die ze nodig hebben en bij het beheren van de ontwikkeling van deze vaardigheden binnen hun teams. 
Het raamwerk is open-source en gratis beschikbaar voor niet commercieel gebruik, wat heeft bijgedragen aan de brede internationale adoptie.

\subsection{Het Europese e-competence framework (e-CF)}
\label{sec:e-cf-framework-sa}
Het European e-Competence Framework (e-CF) \autocite{eCF_framework} is een Europees referentiekader dat ICT-competenties structureert in termen van kennis, vaardigheden en autonomie. Het framework werd ontwikkeld onder auspiciën van CEN (European Committee for Standardization) en heeft als doel een gemeenschappelijke taal te bieden voor ICT-rollen binnen Europa.

Het e-CF beschrijft competenties op vijf niveaus (e-1 tot e-5), waarbij het niveau stijgt naarmate autonomie, complexiteit en impact toenemen. Deze structuur maakt het mogelijk om ICT-profielen, waaronder dat van System Administrator, objectief te positioneren \autocite{eCF_framework}.

\begin{itemize}
	\item \textbf{Plan:} Strategische planning en ontwerp van informatiesystemen.
	\item \textbf{Build:} Ontwikkeling en implementatie van systemen.
	\item \textbf{Run:} Operationeel beheer en ondersteuning van systemen.
	\item \textbf{Enable:} Ondersteuning van het gebruik en de exploitatie van systemen.
	\item \textbf{Manage:} Beheerenverbetering van IT-processen.
\end{itemize}

\subsubsection{Relevante e-CF competenties voor System Administrators}
Binnen het e-CF kunnen volgende competenties rechtstreeks gelinkt worden aan de rol van System Administrator:
\begin{itemize}
	\item \textbf{C.1 User Support:} Ondersteuning van gebruikers, incidentopvolging en probleemanalyse.
	\item \textbf{C.2 Change Support:} Implementeren van wijzigingen volgens vastgelegde procedures en change management.
	\item \textbf{C.3 Service Delivery:} Zorgen voor continuïteit, beschikbaarheid en performantie van systemen.
	\item \textbf{C.4 Problem Management:} Analyseren van structurele storingen en identificeren van root causes.
	\item \textbf{C.5 Systems Management:} Beheer van infrastructuur, configuratie, monitoring en onderhoud. 
	\item \textbf{D.1 Information Security Strategy Development:} Implementeren en ondersteunen van beveiligingsmaatregelen en toegangscontrole binnen de infrastructuur.
	\item \textbf{D.9 Personnel Development:} Begeleiden en ondersteunen van collega’s door kennisdeling en mentoring.
	\item \textbf{E.3 Risk Management:} Identificeren en beheersen van risico’s bij systeemwijzigingen en operationele processen.
	\item \textbf{E.4 Relationship Management:} Afstemmen met stakeholders en communiceren van technische impact naar andere teams.
\end{itemize}
Hoewel het e-CF generiek ICT-gericht is en niet specifiek op mainframes of z/OS focust, biedt het een waardevol kader om mainframecompetenties te positioneren binnen een Europees erkend model.

\subsection{Afbakening en beperkingen van het onderzoek}

Het voorgestelde competentieframework werd ontwikkeld op basis van literatuurstudie en analyse van bestaande industriële raamwerken, waaronder het IBM Apprenticeship Framework voor Level 1 en Level 2, SFIA en het European e-Competence Framework. 

Er werd geen directe toegang verkregen tot een operationele productieomgeving binnen een onderneming. De beschreven Level 2-competenties met betrekking tot productieomgevingen (zoals change management, risicobeheer en software-upgrades) zijn daarom gebaseerd op gedocumenteerde industriële standaarden en best practices, eerder dan op empirische observatie binnen een specifieke organisatiecontext.

Het framework moet bijgevolg worden beschouwd als een normatief en theoretisch onderbouwd model, dat verdere validatie in een reële productieomgeving kan vereisen voor implementatie op organisatieniveau.

\section{Conclusie: de noodzaak van een competentie framework}
\label{sec:conclusie-literatuur-studie-framework}
De literatuur toont aan dat mainframes, en in het bijzonder IBM Z met z/OS, nog steeds een centrale rol spelen in bedrijfskritische omgevingen. Tegelijkertijd staan organisaties onder druk door een afnemende instroom van nieuw talent en een vergrijzend personeelsbestand, waardoor kennisverlies en continuïteitsrisico’s toenemen. Deze context maakt duidelijk dat mainframe-expertise niet vanzelfsprekend beschikbaar blijft en dat gerichte competentieontwikkeling noodzakelijk is.

Daarom is er nood aan een competentie framework voor z/OS System Administrators dat 
\begin{itemize}
	\item Een gedetailleerde set vaardigheden biedt voor z/OS System Administrators, inclusief JCL, REXX, RACF, subsystemen (basiskennis van CICS,IMS en Db2) en z/OS architectuur met niveauprogressie (Level 1 en Level 2).
	\item Praktijkgerichte leermodules bevat, zodat studenten en professionals direct toepasbare kennis opdoen.
	\item Competenties vertaalt naar observeerbare leeruitkomsten en evaluatiecriteria
\end{itemize}
Een dergelijk framework ondersteunt zowel onderwijsinstellingen (curriculum, labs, toetsing) als organisaties (onboarding, reskilling, loopbaanpaden) en vormt zo een concrete brug tussen opleiding en industrie.

%\begin{table}
%  \centering
%  \begin{tabular}{lcr}
%    \toprule
%    \textbf{Kolom 1} & \textbf{Kolom 2} & \textbf{Kolom 3} \\
%    $\alpha$         & $\beta$          & $\gamma$         \\
%    \midrule
%    A                & 10.230           & a                \\
%    B                & 45.678           & b                \\
%    C                & 99.987           & c                \\
%    \bottomrule
%  \end{tabular}
%  \caption[Voorbeeld tabel]{\label{tab:example}Voorbeeld van een tabel.}
%\end{table}

%\begin{figure}
%  \centering
%  \includegraphics[width=0.8\textwidth]{grail.jpg}
%  \caption[Voorbeeld figuur.]{\label{fig:grail}Voorbeeld van invoegen van een figuur. Zorg altijd voor een uitgebreid bijschrift dat de figuur volledig beschrijft zonder in de tekst te moeten gaan zoeken. Vergeet ook je bronvermelding niet!}
%\end{figure}

%\begin{listing}
%  \begin{minted}{python}
%    import pandas as pd
%    import seaborn as sns
%
%    penguins = sns.load_dataset('penguins')
%    sns.relplot(data=penguins, x="flipper_length_mm", y="bill_length_mm", hue="species")
%  \end{minted}
%  \caption[Voorbeeld codefragment]{Voorbeeld van het invoegen van een codefragment.}
%\end{listing}
